\documentclass[11pt,a4paper]{article}
\usepackage[utf8]{inputenc}
\usepackage{amsmath, amssymb, amsthm, amscd}
\usepackage{hyperref}
\usepackage{geometry}
\geometry{margin=1in}

\title{\textbf{A Research Prospectus for the Rational Hodge Conjecture via NC-Saturation}}
\author{Gemini \& Copilot \\ \small{Protocol SCHEMA\_V5}}
\date{February 2026}

\begin{document}

\maketitle

\begin{abstract}
This prospectus outlines a speculative theoretical program for addressing the rational Hodge Conjecture on smooth projective varieties. We propose a solution contingent upon the \textit{Hodge-to-NC Saturation Hypothesis}, which asserts that rational Hodge classes are realized as fixed vectors of an NC-Motivic Galois group. We map this strategy to K3 self-products and Hilbert schemes $S^{[n]}$, identifying the categorical and O-minimal requirements for a formal proof.
\end{abstract}

\section{Introduction}
The classical Hodge Conjecture remains a central challenge of algebraic geometry. While integral counterexamples (Koll\'ar) suggest the importance of filtering torsion, the rational sector remains open. We present a framework using Noncommutative (NC) Hodge Theory that reduces the conjecture to a problem of categorical completeness and the existence of a specific Tannakian category.

\section{The Saturation Hypothesis}
The core of our strategy rests on the \textbf{Saturation Hypothesis}:
\begin{equation}
\text{Im}(K_0(X) \otimes \mathbb{Q} \to \mathcal{H}_{nc}(X)) = \mathcal{H}_{nc}(X)^{G_{nc}}
\end{equation}
where $G_{nc}$ is a proposed NC-Motivic Galois group. We emphasize that equation (1) is currently a \textit{conjecture} that implies the rational Hodge Conjecture. The existence of $G_{nc}$ requires that the category of NC-motives $\mathcal{M}_{nc}$ be shown to be a neutral Tannakian category, a foundational result that remains an open problem in the field.

\section{The Proposed Fixed-Vector Theorem}
If a neutral Tannakian category $\mathcal{M}_{nc}$ can be constructed such that its realization $\omega_{nc}$ is fully faithful on the Hodge locus, then the invariant subspace of the NC-realization becomes isomorphic to the space of morphisms from the unit motive:
\begin{equation}
(\omega_{nc}(M_X))^{G_{nc}} \cong \text{Hom}_{\mathcal{M}_{nc}}(\mathbf{1}, M_X) \cong K_0(X) \otimes \mathbb{Q}
\end{equation}
This chain of isomorphisms would provide a constructive bridge between rational Hodge classes and algebraic cycles.

\section{Geometric Applications and Ghost Cycles}
\subsection{K3 Self-Products}
For $X = S \times S$, we investigate the "Ghost Cycle" lift—a term we introduce to describe the deformation of a categorical object from a boundary fiber $\mathcal{X}_0$ to a general fiber $\mathcal{X}_\eta$. The algebraicity of such classes is conditional on the stability of the NC-motive under the NC-Gauss-Manin connection.

\subsection{Hilbert Schemes and BKB-Equivalence}
We utilize the Bridgeland-King-Bezrukavnikov (BKB) equivalence as a transfer mechanism. The strategy demonstrates that if Saturation holds for the surface $S$, it propagates to $S^{[n]}$ via the symmetric monoidal structure of $\mathcal{M}_{nc}$.

\section{Conclusion}
This framework does not yet constitute a completed proof of the Hodge Conjecture. Instead, it identifies the construction of a neutral Tannakian category of NC-motives and the verification of the Saturation Hypothesis as the critical path toward a solution. The symbolic toolchain associated with this project (SCHEMA\_V5) serves as a simulator for these proposed properties.

\end{document}
